% TOP_PROBLEM: {"id": 30006931, "title": "Exact value of the site-percolation threshold on the square lattice", "category_id": 19, "category": {"id": 19, "name": "probability", "display_name": "Probability", "description": "Problems involving probability theory, stochastic processes, and random structures.", "slug": "probability", "order_index": 19, "created_at": "2026-07-31T15:26:25.671Z"}, "status": "open"}
% FIRST_POSTED: 2026-09-27
% !TeX program = lualatex
% ATTEMPT_STATUS: unresolved
\documentclass[11pt]{article}
\usepackage[margin=25mm]{geometry}
\usepackage{fontspec}
\setmainfont{Latin Modern Roman}
\usepackage{amsmath,amssymb,mathtools}
\usepackage{unicode-math}
\setmathfont{Latin Modern Math}
\usepackage{xurl,hyperref}
\hypersetup{colorlinks=true,urlcolor=blue}
\setcounter{secnumdepth}{0}
\setlength{\parindent}{0pt}
\setlength{\parskip}{5pt}
\setlength{\emergencystretch}{3em}
\begin{document}
\section*{\texorpdfstring{Exact value of the site-percolation threshold on the square lattice}{Exact value of the site-percolation threshold on the square lattice}}\label{30006931}
\subsection{Definitions and mathematical statement}
On the nearest-neighbor graph of ${\mathbb{Z}}^2$, declare each vertex open independently with probability $p$; edges connect neighboring open vertices. Define
\[
 \theta_{\rm site}(p)={\mathbb{P}}_p(0\text{ belongs to an infinite open component}),
 \qquad
 p_c^{\rm site}=\inf\{p\in[0,1]:\theta_{\rm site}(p)>0\}.
\]
Determine this real number exactly. An exact determination requires a rigorously specified value and proof of equality, not a numerical estimate without exact identification. The model is independent site percolation on the square lattice, not bond percolation or site percolation on the triangular lattice.
\par\smallskip\textit{Formulation and concept references:} \href{https://arxiv.org/abs/1507.03027}{[S1]}.

\subsection{Short English statement}
What is the exact critical probability at which independent site percolation on the square lattice can produce an infinite open cluster?
\subsection{Research attempt}
\paragraph{Process and attainable partial result.}
This attempt was generated by GPT-6 Astra Ultra on 27 September 2026. We derive explicit, deliberately coarse rigorous bounds by counting open paths and closed boundaries, then test an exact finite-box calculation as a proposed route to the infinite-lattice threshold. The result is an interval containing the critical probability, not an exact determination or a competitive numerical estimate.

The formulation source [S1] gives a high-precision estimate through transfer-matrix calculations and extrapolation. Its numerical precision must not be read as a proof of an exact algebraic identity. We do not substitute the square-lattice bond threshold or the triangular-lattice site threshold. Those models have different independent variables or different adjacency relations.

\paragraph{Coupling and the order of the quantifiers.}
Assign independent uniform random variables $U_v$ to all vertices, and declare $v$ open at parameter $p$ when $U_v\le p$. This couples all parameters so that every open path at $p$ remains open at any larger parameter. In particular $\theta_{\rm site}(p)$ is nondecreasing.

An infinite component containing the origin has a simple path of every finite length starting there. To see this, a locally finite graph has only finitely many vertices within any fixed graph distance. An infinite connected component therefore contains vertices at arbitrarily large distance, and a shortest path to any such vertex is simple. This implication allows a union bound over finite simple paths to prove nonpercolation. Conversely, positive probabilities of long paths with no bound uniform in their length do not alone prove positive probability of an infinite path.

\paragraph{A path-counting improvement over the first branching bound.}
Let $c_n$ be the number of nearest-neighbor self-avoiding walks of $n$ edges from the origin. The crude estimate is $c_n\le4\cdot3^{n-1}$: four first steps and at most three nonbacktracking choices afterward. We can exclude two more patterns in every group of three subsequent steps.

Fix the directed edge by which a walk just arrived at its current vertex. There are $3^3=27$ possible nonbacktracking three-step extensions before any other exclusions. Two extensions make three consecutive left turns or three consecutive right turns relative to that incoming edge. Together with the incoming edge, each traces a unit square and returns to the preceding vertex. Both are forbidden for a self-avoiding walk, independently of its earlier history. There are therefore at most $25$ legal extensions over the next three steps. Other collisions only reduce this number.

Writing $n-1=3a+r$, $0\le r<3$, gives
\[
 c_n\le4\cdot25^a3^r\le C\mu^n,
 \qquad \mu=25^{1/3}, \tag{433.1}
\]
for a fixed absolute constant $C$. Every simple path of length $n$ uses $n+1$ distinct sites and is open with probability $p^{n+1}$. Thus
\[
 \theta_{\rm site}(p)\le c_np^{n+1}\le Cp(p\mu)^n.
\]
Letting $n\to\infty$ proves $\theta_{\rm site}(p)=0$ for $p<25^{-1/3}$, and hence
\[
 p_c^{\rm site}\ge25^{-1/3}>\frac13. \tag{433.2}
\]
The proof makes no assertion about percolation exactly at the displayed lower endpoint. That is unnecessary for the inequality concerning the infimum.

\paragraph{A closed-boundary argument for a strict upper bound below one.}
Call two vertices star-neighbors if their coordinate differences have maximum absolute value one. Each vertex has eight star-neighbors. The elementary planar boundary lemma used here says that if the origin is open and its nearest-neighbor open component is finite, that component is surrounded by a circuit of distinct closed star-neighbor vertices. One way to obtain it is to follow the exterior boundary of the union of unit squares centered at the vertices of the finite component. Centers of successive exterior neighboring squares coincide or are star-neighbors; exterior squares adjacent across a boundary edge must be closed. Removing repetitions leaves a surrounding closed star circuit. This is the site version of the usual planar boundary argument; diagonal adjacency is needed on the closed side.

For a surrounding circuit with $m$ vertices, its polygonal trace meets the positive horizontal ray at a vertex $(r,0)$ with $1\le r\le m$. The trace crosses that ray at an integer-height endpoint or along a horizontal edge because each step changes the vertical coordinate by at most one. Its horizontal span is at most $m$, which bounds $r$ when the circuit surrounds the origin. For each of these at most $m$ anchors, the number of oriented nonbacktracking star walks with $m$ steps is at most $8\cdot7^{m-1}$. This overcounts all the circuits and is sufficient. Therefore their number is at most $8m7^{m-1}$.

Put $q=1-p$. A specified circuit of $m$ distinct sites is closed with probability $q^m$. A union bound, enlarged by including all $m\ge1$, gives for $7q<1$
\[
 \begin{split}
 \mathbb P_p(0\text{ open with finite component})
 &\le\sum_{m\ge1}8m7^{m-1}q^m\\
 &=\frac{8q}{(1-7q)^2}.
 \end{split} \tag{433.3}
\]
It follows that $\theta_{\rm site}(p)\ge p-8q/(1-7q)^2$. At $q=1/32$, the subtracted expression is $256/625$, strictly less than $31/32$. Consequently
\[
 25^{-1/3}\le p_c^{\rm site}\le\frac{31}{32}. \tag{433.4}
\]
These bounds are intentionally weak. Their value for this attempt is that every estimate has an explicit direction and does not appeal to an extrapolated critical value.

\paragraph{An exact finite crossing polynomial.}
Take the $2$ by $2$ array of vertices with its square-lattice edges, and ask for an open path from its left column to its right column, using only this array. Such a crossing exists exactly when the two sites in the top row are both open or the two sites in the bottom row are both open. The rows use disjoint sites, so the crossing probability is
\[
 C_2(p)=1-(1-p^2)^2=2p^2-p^4. \tag{433.5}
\]
The parameter at which this probability is $1/2$ is exactly
\[
 p_2=\sqrt{1-1/\sqrt2}.
\]
This algebraic number is a property of the stated finite event. Equating it to $p_c^{\rm site}$ would require a theorem linking this particular crossing probability to the infinite-volume transition. No such theorem has been proved in the calculation. Changing the box or the finite event changes the polynomial, so exact arithmetic inside the box does not remove that missing link.

For an arbitrary finite set of $N$ sites, every specified crossing event has a polynomial probability
\[
 C(p)=\sum_{\omega\in\{0,1\}^N}
       1_{\{\omega\text{ crosses}\}}p^{|\omega|}(1-p)^{N-|\omega|}.
 \tag{433.6}
\]
Exhaustive enumeration can certify its coefficients. It cannot by itself prove that a chosen sequence of roots identifies the threshold, or bound the error in that identification at every size. A successful use of larger boxes needs a quantitative finite-size criterion with verified hypotheses and a rigorous connection to the infinite graph.

\paragraph{Why the familiar self-duality shortcut does not apply.}
The complementary boundary in (433.3) uses star-neighbor connections between closed sites, whereas the open model uses four nearest neighbors. Exchanging open and closed therefore exchanges two different adjacency structures. It does not map this site model to itself. This concrete mismatch is already visible around one isolated open vertex: its four closed nearest neighbors form a surrounding star circuit by diagonal steps, but need not be connected by nearest-neighbor closed edges.

Similarly, replacing each edge by the event that its endpoints are open does not create independent bond percolation. Two incident edge events share a site. For distinct $u,v,w$ with edges $uv,vw$, their joint probability is $p^3$, while the product of their separate probabilities is $p^4$. Independence fails for $0<p<1$, so an exact theorem for independent bonds cannot be imported under this replacement.

\paragraph{Verification and outcome.}
Finite enumeration checks (433.1) against all self-avoiding walks through a fixed small length and verifies (433.5) on all sixteen site configurations. Rational arithmetic checks the strict positivity following (433.3). The planar boundary argument and the infinite-length limit are supplied separately; finite enumeration does not certify either limiting statement by itself.

\textbf{UNRESOLVED.} We obtain explicit rigorous bounds and an exact finite crossing calculation, but no exact identification of the infinite-lattice site threshold. The proposed extrapolation route stops at its missing finite-to-infinite theorem.

\subsection{Sources}
\begin{itemize}
\item[S1]Jesper Lykke Jacobsen, Critical points of Potts and O(N) models from eigenvalue identities in periodic Temperley-Lieb algebras, 2015, Section 1(ii) and Eq. (2).
\url{https://arxiv.org/abs/1507.03027}.
\textit{Formulation source.}
\end{itemize}
\end{document}
